% Introduction: draft notes/draft_abstract_intro_2026-10-29.md (inbox 20261029-183827), P2/P4 aligned with the
% revised abstract (inbox 20261029-194129 item 3). Numbers: see notes/table_sources.md 'numbers in text'.
\section{Introduction}\label{sec:intro}
\PARstart{A}{utonomous} vehicles and mobile robots carry several sensors because each one fails in different conditions. A camera loses contrast at night and in fog, a depth estimate develops holes on reflective surfaces, an event camera goes silent on static scenes, and LiDAR returns thin out in rain and with distance. Multimodal semantic segmentation fuses these sensors so that the scene is still understood when one of them degrades.

The fusion models that lead the full-sensor benchmarks~\cite{cmnext2023,cmx2023,stitchfusion2024,geminifusion2024,cafuser2025,dgfusion2026} handle adverse conditions by telling the fusion what the condition is. CAFuser~\cite{cafuser2025} infers a condition token from the camera image and supervises it with a textual description of the weather and lighting, and DGFusion~\cite{dgfusion2026} adds local depth tokens and depth supervision. These models see the adverse scenes contained in their training data, but they are not trained with degraded or missing sensors, and they lose much of their accuracy when a sensor actually fails: under one evaluation protocol on DELIVER, the official DGFusion weights fall from 56.7 to 31.1 mean intersection over union (mIoU) when averaged over all subsets of removed sensors. Methods designed for missing sensors~\cite{magic2024,anyseg2024,rmmss2025,fuseadapt2026} train with modality dropout and distillation, but the two lines of work rarely evaluate the same models under the same protocol~\cite{mmssbench2025}. Frozen vision foundation models offer strong features for every sensor, yet their multimodal adaptations mostly keep RGB as the primary stream and inject one auxiliary sensor~\cite{mmsamadapter2025}.

We ask whether a change of the training signal, instead of additional condition supervision, can give an accurate fusion model the ability to hold up when its sensors degrade. We propose Clean-to-Degraded Two-Pass Distillation (C2D2). Every training batch is used twice, once with clean inputs and once with sensor-specific synthetic degradation, such as blur and exposure changes for the camera, holes for depth, dropped beams for LiDAR, and lost resolution for events, together with entire sensor removal. The degraded pass is supervised by the labels and distilled from a frozen teacher that sees the clean input. C2D2 needs no condition or degradation labels and adds no computation at inference. We apply it to a fusion model that reads every sensor with one frozen DINOv3 backbone~\cite{dinov3_2025} through sensor-specific low-rank adapters, fuses the sensors with cross-modal attention, and recovers thin structures with a high-resolution decoding path.

On DELIVER~\cite{cmnext2023}, the model reaches 57.74 mIoU on the test set (mean of three training runs), 1.0 point above the official weights of DGFusion, the best published four-sensor fusion model, evaluated under the same protocol, and it keeps 48.7 mIoU when averaged over all subsets of removed sensors, compared with 31.1 for the official DGFusion weights and at most 46.5 for strong fusion models retrained with the same degradation \todo{MUSES and MCubeS results (FU-08)}. We also analyse how the dependence on each sensor changes across conditions, where the high-resolution path helps, and where the model still fails.

Our contributions are as follows.
\begin{itemize}
\item We propose C2D2, a label-free two-pass training scheme that combines sensor-specific synthetic degradation with distillation from a clean-input teacher and adds no inference cost. \todo{after FU-03 (L7) and FU-04 (DGFusion/CAFuser transfer): one sentence on how it differs from missing-modality distillation and whether its gain depends on the frozen foundation backbone}
\item We build a single fusion model on a frozen vision foundation backbone that, trained with C2D2, improves full-sensor accuracy on DELIVER \todo{and MUSES and MCubeS (FU-08)} while remaining robust under the public missing-sensor protocol.
\item We analyse sensor dependence across conditions, missing, partial, and noisy sensors, thin objects, and failure cases, and report where the model still fails. \todo{state the code and checkpoint release}
\end{itemize}
